% GENERATED FILE. Edit canonical lessons, then run npm run generate:books.
% canonical-source-hash: fnv1a64:2fa07f26f4876709
% canonical-lessons: LA-C276-vespa, LA-C276-culex, LA-C276-pulex, LA-C276-vermis, LA-C276-testudo

\chapter{Wasp, Gnat, Flea, Worm, Tortoise}
\label{ch:la-wasp-gnat-flea-worm-tortoise}
\hlchaptermodality{276}

\begin{chapteropening}
\textbf{By the end of this chapter:} \emph{I can say a wasp, a gnat, a flea, a worm and a tortoise.}

Complete the last lesson of chapter 276: I can say a wasp, a gnat, a flea, a worm and a tortoise.
\end{chapteropening}

\section[vespa, vespae]{vespa, vespae — a wasp}

\label{lesson:LA-C276-vespa}



\begin{quote}
Before the new one: say the Latin for smoke, then the Latin for ash.
\end{quote}

\subsection*{You'll want to know: vespa, vespae}
\textbf{vespa, vespae} — \textquotedblleft{}a wasp\textquotedblright{}.

\begin{cousinweb}
\textbf{vespa} is inherited, usually traced to Proto-Indo-European \emph{*wobhseh2}, \textquotedblleft{}wasp\textquotedblright{}, from \emph{*webh-}, \textquotedblleft{}to weave\textquotedblright{} (EtymDB 2.1): the insect that weaves its nest. English \emph{wasp} is a cousin (Lewis \& Short).
\end{cousinweb}

\begin{grammarlens}[title={What you've built}]
One more word for this chapter.
\end{grammarlens}

\subsection*{Your turn}
\begin{itemize}
\raggedright
  \item \emph{Say it:} \emph{vespa, vespae}
  \item \emph{Say it:} \emph{vespa, vespae}, once more
  \item \emph{Say it:} say \emph{vespa, vespae}
  \item \emph{Recall:} say the Latin for smoke, then the Latin for ash, then say \emph{vespa, vespae} again
\end{itemize}

\begin{tcolorbox}[breakable,colback=teal!4,colframe=teal!35!black,title={Before you move on}]
What does vespa, vespae mean? (\textquotedblleft{}A wasp\textquotedblright{}.) Say it once more.
\end{tcolorbox}
\section[culex, culicis]{culex, culicis — a gnat, a mosquito}

\label{lesson:LA-C276-culex}



\begin{quote}
Before the new one: say the Latin for ash, then the Latin for a wasp.
\end{quote}

\subsection*{You'll want to know: culex, culicis}
\textbf{culex, culicis} — \textquotedblleft{}a gnat, a mosquito\textquotedblright{}.

\begin{cousinweb}
\textbf{culex} has no sure origin; EtymDB 2.1 records a proposed Proto-Indo-European source, but the link is not settled.
\end{cousinweb}

\begin{grammarlens}[title={What you've built}]
One more word for this chapter.
\end{grammarlens}

\subsection*{Your turn}
\begin{itemize}
\raggedright
  \item \emph{Say it:} \emph{culex, culicis}
  \item \emph{Say it:} \emph{culex, culicis}, once more
  \item \emph{Say it:} say \emph{culex, culicis}
  \item \emph{Recall:} say the Latin for ash, then the Latin for a wasp, then say \emph{culex, culicis} again
\end{itemize}

\begin{tcolorbox}[breakable,colback=teal!4,colframe=teal!35!black,title={Before you move on}]
What does culex, culicis mean? (\textquotedblleft{}A gnat, a mosquito\textquotedblright{}.) Say it once more.
\end{tcolorbox}
\section[pūlex, pūlicis]{pūlex, pūlicis — a flea}

\label{lesson:LA-C276-pulex}



\begin{quote}
Before the new one: say the Latin for a wasp, then the Latin for a gnat.
\end{quote}

\subsection*{You'll want to know: pūlex, pūlicis}
\textbf{pūlex, pūlicis} — \textquotedblleft{}a flea\textquotedblright{}.

\begin{cousinweb}
\textbf{pūlex} is inherited, usually traced to Proto-Indo-European \emph{*plusis}, \textquotedblleft{}flea\textquotedblright{} (EtymDB 2.1), with its sounds rearranged; Greek \emph{psylla}, \textquotedblleft{}flea\textquotedblright{}, is a cousin.
\end{cousinweb}

\begin{grammarlens}[title={What you've built}]
One more word for this chapter.
\end{grammarlens}

\subsection*{Your turn}
\begin{itemize}
\raggedright
  \item \emph{Say it:} \emph{pūlex, pūlicis}
  \item \emph{Say it:} \emph{pūlex, pūlicis}, once more
  \item \emph{Say it:} say \emph{pūlex, pūlicis}
  \item \emph{Recall:} say the Latin for a wasp, then the Latin for a gnat, then say \emph{pūlex, pūlicis} again
\end{itemize}

\begin{tcolorbox}[breakable,colback=teal!4,colframe=teal!35!black,title={Before you move on}]
What does pūlex, pūlicis mean? (\textquotedblleft{}A flea\textquotedblright{}.) Say it once more.
\end{tcolorbox}
\section[vermis, vermis]{vermis, vermis — a worm}

\label{lesson:LA-C276-vermis}



\begin{quote}
Before the new one: say the Latin for a gnat, then the Latin for a flea.
\end{quote}

\subsection*{You'll want to know: vermis, vermis}
\textbf{vermis, vermis} — \textquotedblleft{}a worm\textquotedblright{}.

\begin{cousinweb}
\textbf{vermis} is inherited, usually traced to Proto-Indo-European \emph{*wrmis}, \textquotedblleft{}worm\textquotedblright{} (EtymDB 2.1); Gothic \emph{waurms}, \textquotedblleft{}serpent\textquotedblright{}, and English \emph{worm} are cousins (Lewis \& Short).
\end{cousinweb}

\begin{grammarlens}[title={What you've built}]
One more word for this chapter.
\end{grammarlens}

\subsection*{Your turn}
\begin{itemize}
\raggedright
  \item \emph{Say it:} \emph{vermis, vermis}
  \item \emph{Say it:} \emph{vermis, vermis}, once more
  \item \emph{Say it:} say \emph{vermis, vermis}
  \item \emph{Recall:} say the Latin for a gnat, then the Latin for a flea, then say \emph{vermis, vermis} again
\end{itemize}

\begin{tcolorbox}[breakable,colback=teal!4,colframe=teal!35!black,title={Before you move on}]
What does vermis, vermis mean? (\textquotedblleft{}A worm\textquotedblright{}.) Say it once more.
\end{tcolorbox}
\section[testūdō, testūdinis]{testūdō, testūdinis — a tortoise}

\label{lesson:LA-C276-testudo}



\begin{quote}
Before the new one: say the Latin for a flea, then the Latin for a worm.
\end{quote}

\subsection*{You'll want to know: testūdō, testūdinis}
\textbf{testūdō, testūdinis} — \textquotedblleft{}a tortoise\textquotedblright{}.

\begin{cousinweb}
\textbf{testūdō} is built on \emph{testa}, \textquotedblleft{}a piece of baked clay, a shell\textquotedblright{} (Lewis \& Short), with \emph{-ūdō}, an ending for a quality: the shelled creature.
\end{cousinweb}

\begin{grammarlens}[title={What you've built}]
One more word for this chapter.
\end{grammarlens}

\subsection*{Your turn}
\begin{itemize}
\raggedright
  \item \emph{Say it:} \emph{testūdō, testūdinis}
  \item \emph{Say it:} \emph{testūdō, testūdinis}, once more
  \item \emph{Say it:} say \emph{testūdō, testūdinis}
  \item \emph{Recall:} say the Latin for a flea, then the Latin for a worm, then say \emph{testūdō, testūdinis} again
\end{itemize}

\begin{tcolorbox}[breakable,colback=teal!4,colframe=teal!35!black,title={Before you move on}]
What does testūdō, testūdinis mean? (\textquotedblleft{}A tortoise\textquotedblright{}.) Say it once more.
\end{tcolorbox}
